\documentclass[12pt,a4paper,oneside]{article}
\usepackage[brazil]{babel}
\usepackage[utf8]{inputenc}
\usepackage{olymp}
\usepackage{color}
\usepackage[dvipsnames]{xcolor}
\usepackage{listings}

\setcounter{problem}{0}

\begin{document}

\begin{problem}{Combinação}{combinacao}{2}
 
Uma turma de faculdade tem $N$ alunos e precisa escolher $K$ para compor a comissão de organização da calourada. O número de diferentes comissões que
podem ser formadas pode ser calculado pela fórmula do coeficiente binomial, que consiste no número de combinações de $N$ termos, $K$ a $K$:

\begin{center}
$\displaystyle \binom{N}{K} = \frac{N!}{K!(N-K)!}$
\end{center}

Como a fórmula requer o cálculo de fatoriais, foi criada a função abaixo:

\lstset{language=C++,
                basicstyle=\ttfamily,
                keywordstyle=\color{Mulberry}\ttfamily,
                stringstyle=\color{Maroon}\ttfamily,
                commentstyle=\color{YellowGreen}\ttfamily,
                morecomment=[l][\color{magenta}]{\#},
                basewidth = {.48em}
}
\begin{lstlisting}
//Retorna o valor de x!
long long int fatorial (int x) {
   long long int fat = 1;
   
   for(int i=2; i<=x; i++)
      fat *= i;
      
   return fat;
}
\end{lstlisting}

Sua tarefa é fazer um programa que utiliza esta função para calcular o número de diferentes comissões.

%\Notes
%
%Observações, se tiver.

\InputFile

A entrada contém dois valores inteiros, $N$ e $K$, que são respectivamente o número de alunos da turma e o número de alunos na comissão. Restrição: $1 \leq K \leq N \leq 20$.


\OutputFile

Seu programa deve gerar apenas uma linha de saída, o número de
diferentes comissões que podem ser formadas.

\Notes
Nesta questão, seu programa deve usar a função \texttt{fatorial} dada acima exatamente como está, não a modifique! Sua
tarefa é apenas criar o programa com a função \texttt{main}.

\Example

\begin{example}
\exmp{
12 4
}{
495
}%
\end{example}

\begin{example}
\exmp{
7 7
}{
1
}%
\end{example}

\begin{example}
\exmp{
15 5
}{
3003
}%
\end{example}


%Comentários sobre o exemplo, se necessário.

\end{problem}

\end{document}
